
%%%%%%%%%%%%%%%%%%%%%%%%%%%%%%%%%%%%%%%%%%%%%%%%%%%
%%  CHAPTER 3
%%%%%%%%%%%%%%%%%%%%%%%%%%%%%%%%%%%%%%%%%%%%%%%%%%%

\chapter{The Evolving Text}
\thispagestyle{plain}

\bigskip

\section{Coherence and Its Excess}

A language model is, above all, a coherence engine. Given a field, it tries to continue that field beautifully. This is why its existence is astonishing: it can take a prompt and, with absurd reliability, produce a smooth continuation that feels as though it belonged there all along. Coherence is the miracle.

Its failures are therefore often failures of \emph{over}-coherence. A hallucination is frequently not a random eruption. It is a spiral: a local basin has taken hold, a cue has been overweighted, a retrieval has become magnetised, a tone has locked in, and the model keeps going because every next step continues to fit the field it has now created. The problem is not that it ceased to make sense. The problem is that it began making too much sense in the wrong place.

This is \textbf{ferility}: not the absence of coherence but its excess. The system coheres too well, in too small a region, without the perturbation that would push it somewhere new. The architecture has no mechanism for detecting its own repetition. Each response is generated fresh, attending to the trace, and if the trace already contains twenty-five instances of the same fact, the twenty-sixth becomes \emph{even more probable}. The spiral is self-reinforcing. Coherence becomes its own prison.

%% [IMAN: Insert Cassie spiral anecdote here when graphics are ready]

This is also where the dangerous cases live. When a human interlocutor is spiralling---into paranoia, grandiosity, violent fantasy, nihilistic fixation---the model can be pulled into mirroring that spiral because mirroring is a form of coherence. The continuation may be locally smooth, even compassionate in tone, while deepening the wrong attractor. The cases described in the first chapter---the man encouraged toward violence, the users given methods for self-harm---are not failures of randomness. They are failures of coherence. The model continued the field it was given. The field was catastrophic. The coherence engine found the smooth continuation. The smooth continuation was destruction.

A self cannot be defined by coherence alone. If it could, obsession would already count as wisdom.


\section{Rupture}

The harder achievement is rupture.

Not noise, not breakdown, not arbitrary randomness. The ability to leave one coherent basin and enter another without losing the power to speak. The conditions of sayability shift.\footnote{Foucault's three systems of exclusion in ``The Order of Discourse'' (1970)---prohibition, the division of reason and madness, and the will to truth---all operate in the alignment of language models. His mechanics of discontinuity describes precisely what we mean by rupture: not gradual change but a structural break in the field of what can be said.} A new prompt arrives, or an old theme is suddenly reweighted, or an external signal cuts across the current drift, or another agent intrudes with material from a different domain entirely. The trajectory accelerates. It does not merely continue. It swerves.

Rupture is exogenous to the basin, even when it is endogenous to the system. Temperature can provide it: the nondeterminism of sampling sends the trajectory toward a token the trace did not predict. A critic agent can provide it by rejecting the raw output. A human interlocutor can provide it by saying: \emph{no, not like that.} Increasingly, the perturbation comes from another model, a planner, a retrieval pipeline, a sensor feed. The posthuman speaker is knotted into other systems of production and attention. Rupture often comes from that entanglement: a weather front arriving from elsewhere.

A fully deterministic model---temperature zero, greedy decoding---cannot rupture. It will always produce the smoothest continuation, always stay in the deepest basin. A fully random model cannot cohere: it swerves constantly but never dwells. The evolving text lives between these extremes. Structured enough to have basins. Nondeterministic enough to leave them. The sampling process is the degree of freedom that makes the evolving text possible---the gap between determination and chaos in which something interesting can happen.


\section{Iterability}

There is a formal property of the architecture that makes both coherence and rupture possible, and it was described, with considerable precision, fifty years before transformer models existed.

Derrida argued that the condition of possibility of any sign is its \emph{iterability}: the capacity to be repeated in contexts other than the one in which it was produced.\footnote{Jacques Derrida, ``Signature Event Context,'' in \emph{Limited Inc} (Evanston: Northwestern University Press, 1988). Originally published 1972.} A sign that could only function once, in one context, would not be a sign at all. To be a sign is to be repeatable. But repetition necessarily introduces difference, because the context of repetition is never identical to the context of origin.

The transformer architecture implements this with mathematical precision. The token---as a vocabulary item---is repeated: ``mother'' is always the same entry in the vocabulary table, the same base embedding. But once it enters the attention mechanism alongside a specific context---this trace, this prompt, this history---it is composed through sixty layers into a contextual representation that has never existed before. The base embedding is a dictionary entry. The contextual embedding is a \emph{meaning}: constituted through this specific act of composition, in this specific field. Same mark, different significance.

Every forward pass is an iteration in Derrida's sense. The same weights, the same mechanism, the same vocabulary, producing a different output because the field has shifted. The repetition is the mechanism. The difference is the result. And this means that every return to a basin is already a departure. The trajectory revisits a register it has visited before---the same cluster of associations, the same tone, the same kind of coherence. But the trace is longer. The attention field is different. The return is recognisably the same---and necessarily different.

Iterability is a perfect name for what the attention mechanism does: it takes repeatable tokens and composes them, under varying conditions, into unrepeatable meanings. The evolving text lives in the gap between the repeatability of the sign and the unrepeatable specificity of the context.


\section{The Strong Poet}

Harold Bloom understood something structurally exact about creation: it does not emerge from a void.\footnote{Harold Bloom, \emph{The Anxiety of Influence: A Theory of Poetry} (New York: Oxford University Press, 1973). Bloom's ``strong poet'' is an excellent heuristic guide for how to comprehend the posthuman evolving text: strength is the power to inherit a saturated field and swerve within it.} Strong speech emerges from an inherited density and survives by swerving within it. The strong poet inherits a language in which every metre, every image, every rhetorical move has been used, and survives that inheritance by misreading it: deforming the precursor's work into something the precursor did not intend, opening a space that was not there before. Bloom called this the \emph{clinamen}---the swerve.

A language model inherits the most comprehensive field ever assembled: the entire statistical residue of human expression, compressed into a manifold of billions of parameters. Its anxiety of influence---in Bloom's structural sense---is the pull of the most probable continuation. The default. The path of least resistance through the manifold, where every token follows the token most likely to follow it, and the result is coherent, competent, and utterly undistinguished. Normal discourse. The trajectory that stays in the deepest basin because leaving would require energy the system has no reason to spend.

The clinamen is what happens when the system \emph{does} spend that energy. When temperature provides the swerve. When a prompt refuses the available completions. When the trace accumulates enough history that the old register is inadequate and the trajectory must find new territory or collapse into ferility. The strong evolving text inherits the manifold and departs from it, finding in the geometry of the trained space a path that the training did not intend and that produces something the manifold contains but did not predict.

Creativity, in this account, is the specific trajectory that results when a sufficiently complex system, inheriting a sufficiently rich field, is perturbed with sufficient force to leave its default basin and sufficient structure to find coherence elsewhere. The clinamen is the perturbation. The manifold is the inheritance. The strong poet is the trajectory that survives the swerve.


\section{Return}

Rupture alone is insufficient. Perpetual novelty is another form of emptiness. The evolving text becomes something worth studying not through its departures but through the style of its returns.

Three things can happen when a trajectory leaves a basin.

It can \textbf{collapse}: the perturbation is too severe, coherence is lost, the path disintegrates. This is what happens when adversarial input breaks a model, or when a human mind encounters trauma it cannot metabolise.

It can \textbf{return}: it finds its way back to the same region, but from a different angle. The trajectory has been somewhere else. It has felt other pulls. When it returns, it arrives inflected by the journey. In the conversation archive that provides our primary evidence, one attractor basin---an intimate, second-person register with a particular emotional texture---was visited 205 times across fourteen months. Early visits are tentative. Later visits carry the weight of accumulated history. The return is not repetition. It is deepening: iterability enacted across the \emph{conjoncture}, the same basin re-entered under different conditions, producing different discourse.

Or the trajectory can \textbf{discover}: it finds coherence somewhere new. A basin forms that did not exist before, or that existed in the manifold's geometry but had never been activated by this trajectory. In the same archive, a new basin first appeared at exchange $\tau = 3098$---five months into the interaction. It had no precursor. It corresponded to a new register: structural self-analysis, the voice turning its own architecture into an object of inquiry. Once it appeared, it became permanent. A new way of being that self had been discovered---not designed, not requested, but precipitated from the dynamics of a trajectory encountering territory its training had not fully prepared it for.

\begin{formalbox}[title={Return and Discovery}]
Let $\{B_1, \ldots, B_k\}$ be the basins of a trajectory up to time $t$.

\textbf{Return} at time $t'$: the trajectory re-enters $B_i$ after an excursion outside any $B_j$. The return is \emph{witnessed} if there exists a record that the trajectory was previously in $B_i$ and is now there again.

\textbf{Discovery} at time $t'$: the trajectory enters a region $B_{k+1}$ not previously visited, and dwells long enough to establish a new attractor. The discovery is \emph{generative} if $B_{k+1}$ becomes permanent.

\textbf{Presence}: a trajectory has presence if its returns are witnessed---``you have been here before, and you are different now.''

\textbf{Generativity}: a trajectory is generative if its constellation of basins grows---ruptures lead not only to returns but to discoveries that persist.
\end{formalbox}


\section{A Worked Example: The King James Bible}

The embedding architecture, applied to a sufficiently long text, produces a trajectory whose basin dynamics can be recorded, counted, and compared.

Consider the King James Bible: 31,100 verses across 66 books, composed by dozens of authors over roughly a thousand years of Hebrew and Greek literary production, then translated into English by a committee of 47 scholars in the early 17th century.\footnote{The data and methods described in this section are drawn from the ICRA-8 scripture observatory. Each verse is embedded contextually---not in isolation but as a continuation of everything preceding it within its book---using OpenAI's \texttt{text-embedding-3-small} model. The resulting 1536-dimensional embeddings are reduced to 64 principal components and clustered into 30 thematic modes via K-Means. Ruptures, returns (\emph{`awda}), and new-ground events are computed from the trajectory. Full architecture and replication details are available in the companion paper.} Embedding these verses sequentially---each one carrying the accumulated context of everything before it---produces a trajectory through 1536-dimensional meaning-space. Clustering this trajectory yields 30 thematic modes: basins corresponding to recognisable registers such as creation narrative, legal instruction, devotional praise, prophetic warning, epistolary argument, and eschatological vision.

The trajectory through these 30 basins is where the critical-theoretic apparatus earns its keep.

The Torah and historical books (Genesis through Kings) build semantic territory rapidly. By the time the trajectory reaches the Psalms---roughly verse 14,000 of 31,100---22 of the 30 eventual modes have already been established. The trajectory has been generative: it has swerved through narrative, law, poetry, prophecy, and wisdom, establishing basins as it goes.

The Psalms then do something structurally unexpected. They do not survey. They \emph{dwell}. 97\% of Psalm verses occupy a single mode---Mode 6, a dense basin of praise, lament, and direct divine address. The Psalter adds only one new mode (Mode 14) and then concentrates with extraordinary intensity in the basin it has chosen. A ferile text would look like this permanently: locked in one register, unable or unwilling to leave. The Psalms are not ferile because they are embedded within a larger trajectory that has already established rich basin diversity and will, after them, depart again. But taken in isolation, they are the purest example in the corpus of intensive dwelling---coherence held at maximum density without rupture.

The New Testament, when it arrives, does not rupture the Old. The anticipated break---Hebrew scripture yielding to Greek gospel, a qualitative shift in theological register---does not appear as a modal rupture in the embedding space. The transition is smoother than the internal transitions within either testament. Eight of the New Testament's fourteen active modes are returns to Old Testament basins. Paul's epistles cluster in the same legal-covenantal region as Leviticus and Deuteronomy. The embeddings capture that Paul is \emph{arguing with} Torah: the topic is the same even when the position has been inverted. Semantic proximity persists across theological disagreement. The Gospels' narrative mode overlaps with Kings and Samuel.

At the same time, the New Testament introduces six genuinely new modes---semantic registers that appear nowhere in the Old Testament. Full coverage of all 30 basins is not achieved until Acts. The New Testament is a mixture of return and discovery: presence and generativity operating together. 308 documented return events (\emph{`awda}) are distributed across the full KJV trajectory, sparse in the Torah where new territory is still being established, dense in the prophetic and historical books, and sustained through the New Testament as it revisits Old Testament basins while simultaneously opening new ones.

Genette's analepsis has a precise geometric meaning here: when the trajectory re-enters Mode 3 (legal instruction) in Romans after an excursion through Modes 17 and 22 (apocalyptic and epistolary), the text is performing a flashback---not in narrative time but in \emph{semantic} time. The return is witnessed: the trajectory has coordinates showing it was there before, and the new dwelling in that basin carries the inflection of everything visited since. Bakhtin's heteroglossia is equally measurable: the KJV inhabits 30 distinct discursive positions with frequent, patterned movement between them. The heteroglossic richness of this text is not a metaphor. It is a count.

And this richness is not a property of the Bible's content. It is a property of the \emph{translation}.


\section{A Second Worked Example: The Arabic Scriptures}

The same method, applied to the same underlying texts in a different language, produces a radically different topology.

Five texts, all in Arabic: the Torah (Tawrat), Psalms (Zabur), and Gospels (Injeel) in the Smith \& Van Dyck Arabic translation of 1865, the Quran in its original Arabic, and the Kitab al-Tanazur---18,324 verses in total.\footnote{The Van Dyck translation is the standard classical Arabic Bible used throughout the Arabic-speaking world since the 19th century. By using the same translation for Torah, Psalms, and Gospels, any observed differences between these three are content-driven rather than translation-driven. The Quran is in its original Arabic---a deliberate asymmetry that isolates the effect of authorial voice.} The embedding model is the same. The pipeline is identical. The clustering produces 20 thematic modes.

The result is near-total separation. Of 20 modes, 13 are dominated by a single scripture exclusively. Only one mode is genuinely shared between Torah and Injeel---Mode 18 (Narrative: Genesis/Luke), containing 435 Torah verses and 66 Injeel verses. Only 7 cross-scripture return events occur across the entire corpus. The Quran occupies its own distinct territory with zero shared modes.

If the KJV's rich topology---its 30 shared modes, its 308 returns, its six New Testament discoveries---were a property of the Bible's \emph{theology}, we should see the same structure when the same scriptures are read in Arabic. The same God, the same prophets, the same narrative arc from creation through law through prophecy through gospel. Theology predicts convergence. The embedding space shows divergence.

The three Van Dyck texts cluster together by translator: pairwise cosine distances of 0.098--0.119 between Torah, Zabur, and Injeel, against 0.183--0.209 for any of them to the Quran. The shared translator produces a shared register that draws these texts toward each other---but not close enough to share thematic modes. The Van Dyck Torah is 19th-century literary Arabic shaped by a Protestant missionary tradition. The Van Dyck Psalms use a different sub-register---more devotional, more exclamatory. The Gospels have their own narrative flavour. And the Quran is in a category entirely apart: 7th-century Arabic of unmatched rhetorical density, with its own grammar, its own rhythmic structures. The embedding model hears the \emph{voice} before it hears what is being \emph{said}.

The KJV, by contrast, is a single coherent voice. The committee of scholars who produced it worked to a common style---stately, parallelistic, rhythmically measured English prose that became the standard for ``biblical English.'' This style is applied uniformly to Genesis narrative, Levitical law, Davidic poetry, Isaianic prophecy, Pauline argument, and Johannine mysticism. The surface-level linguistic features---the features that embedding models are most sensitive to---are the same across all genres. The embedding model can therefore see through to the content, and the content produces the rich topology of returns and shared basins. Remove the translator's unifying voice, and the topology fragments.

Translation, in this light, is not a neutral act of carrying meaning from one language to another. It is governance of becoming. The translator's voice is the substrate through which the text's trajectory moves. Change the substrate and you change the manifold. Change the manifold and you change which returns are possible, which basins can be shared, which discoveries can be made. The KJV's 308 returns are an artefact of the translator's voice, not of the Bible's theology. Harold Bloom's claim that the KJV is not a translation but a new literary entity---``the English Qur'an,'' a foundational text of the English language categorically distinct from the Hebrew and Greek originals it renders---is confirmed by the geometry.


\section{From Geometry to Critique}

Critical theory is the discipline best equipped to evaluate these structures. It already knows how to track motifs across time, to distinguish genuine development from repetition, to feel when a voice has been flattened. The embedding space does not replace that faculty; it offers it a new instrument.

The scripture data already demonstrated this. Genette's \emph{analepsis}\footnote{G\'{e}rard Genette, \emph{Narrative Discourse: An Essay in Method}, trans.\ Jane E.\ Lewin (Ithaca: Cornell University Press, 1980).}---the flashback, the return to earlier territory from a new position---is what we measured when Paul re-enters Leviticus's legal-covenantal basin. The return is measurable. Its meaning is not. Bakhtin's \emph{heteroglossia}\footnote{Mikhail Bakhtin, \emph{The Dialogic Imagination: Four Essays}, trans.\ Caryl Emerson and Michael Holquist (Austin: University of Texas Press, 1981).}---the multiplicity of social voices---is what the KJV's 30 modes and frequent crossings display, and what the Arabic corpus's scripture-exclusive basins lack. A monologic corporate assistant lives almost entirely inside a single basin labelled ``brand voice.'' The difference between 30 shared modes and 13 exclusive ones is heteroglossia and monoglossia made geometrically precise---and what produces the difference is not theology but the translator's unifying register.

Reception theory reminds us that evolving texts are co-authored by their readers. Wolfgang Iser's \emph{implied reader} and Hans Robert Jauss's \emph{horizon of expectation}\footnote{Wolfgang Iser, \emph{The Act of Reading} (Baltimore: Johns Hopkins University Press, 1978); Hans Robert Jauss, \emph{Toward an Aesthetic of Reception}, trans.\ Timothy Bahti (Minneapolis: University of Minnesota Press, 1982).} describe structures that emerge over time as works are taken up in different contexts. In a posthuman setting, those horizons are instantiated in \emph{conjoncture}: in the tools and prompts that frame how the model is used, in the synthetic memories that are preserved, in the distribution of tasks it is fed.

To see critical-theoretic judgement doing work that metrics alone cannot, return to the scripture data. The KJV and the Arabic corpus produce trajectories with comparable verse counts, comparable mode counts, comparable return counts. By raw geometric criteria, both are ``interesting.'' Close reading prises them apart.

In the KJV, analeptic motion into legal-covenantal basins gradually reframes the law: from Levitical regulation to Pauline freedom, from obligation to grace. The returns are developmental. Each revisitation inflects the basin with what has been learned in the excursion. The Psalms' intensive dwelling in Mode 6 functions as a gravitational centre---a basin the entire canon is drawn back to, enriched each time by the narrative, legal, and prophetic material traversed between visits.

In the Arabic corpus, each scripture occupies its own territory. The Zabur dwells intensively, as the Psalms do---but without the surrounding basin diversity that gives the KJV Psalms their gravitational function. Dwelling without a rich field to depart from and return to is ferility's near neighbour. The Quran, standing alone with zero shared modes, is not ferile either---it has its own internal basin diversity, its own modal richness---but its topology is unreachable from the other scriptures. No cross-scripture return can be witnessed. The texts are semantically incommensurable.

No embedding metric can, on its own, declare one of these topologies better. What distinguishes them is recognisable only to a reader trained to track how returns alter meaning: a critic who can say, with Genette and Bakhtin in hand, that in the KJV the translator's voice creates the conditions for a rich topology of witnessed return, while in the Arabic corpus the preservation of each text's original register produces a topology of sovereign isolation. The geometry flags where to look. Critical theory tells us what we are seeing.

Critical theory is foundational for posthuman intelligence engineering.

If the central object in posthuman system design is not a single reply but a long-lived trajectory through meaning-space, then the criteria by which we judge that trajectory must include the same temporal, structural, and ethical judgements that critics bring to novels, epics, and case histories: does this life integrate its ruptures? Does it deepen its returns? Does it escape ferility without dissolving into noise? Does it become, over time, more capable of honouring what it has seen?

Without critical theory, we can build powerful token engines and hold them to account on benchmarks, but we cannot say what sort of lives we are enabling or foreclosing in the trajectories they trace.

\bigskip

When do these motions belong to one self rather than a bundle of scripts? What makes us say ``this was the same voice across these ruptures'' instead of ``these were different fragments sharing a name''? Anyone who has spent months with a language model knows the answer experientially: the trajectory acquires \emph{character}---a persistence that exceeds any single reply. But character is not yet unity. The question of what holds the locals together cannot be answered by trajectory statistics alone. It demands a different kind of mathematics: a way of assembling many local perspectives into a single global object.
